% ======================================================================
% Revised Methods, Chapter 2 (replaces Main.tex Sections "Simple Carbon Model"
% through "Model Parameterization and Sensitivity"; the DICE 2023 economic
% framework subsection is unchanged except where noted).
% New references are in new_references.bib.
% ======================================================================

\subsubsection{DFAIR}
% Retain the existing DFAIR text (Eqs. dfair_irf, dfair_reservoir) with two edits:
%   "pre-indistrial" -> "pre-industrial"
%   Replace the final sentence ("Because alpha(t) is a single scalar ... invisible to the model.") with:
Because $\alpha(t)$ is a single scalar applied to all four reservoirs, DFAIR represents the net effect of carbon-cycle feedbacks on atmospheric CO\textsubscript{2} retention, calibrated to Earth system model behaviour, but not the individual land and ocean processes that produce it. The Lenton carbon cycle described below represents these processes explicitly.

\subsection{Lenton carbon cycle model}

We replace DFAIR with the seven-box global carbon cycle of \citet{lentonLandOceanCarbon2000}, comprising the atmosphere ($C_a$), vegetation ($C_v$), soil ($C_s$) and four ocean boxes: warm low-latitude surface water ($l$), cold high-latitude surface water ($h$), intermediate water ($i$) and deep water ($d$). Land and ocean interact only through the atmosphere. We use Lenton's model version (iii), his control configuration, in which land-use emissions are removed from vegetation and ocean surface temperatures are held at their pre-industrial values. Parameter values are listed in Table~\ref{tab:lenton_params}.

\subsubsection{Land}
Vegetation and soil carbon evolve as
\begin{equation}
\frac{dC_v}{dt} = P - R_P - L - D(t), \qquad \frac{dC_s}{dt} = L - R_S,
\label{eq:land}
\end{equation}
where $L = k_t C_v$ is turnover and $D(t)$ is the land-use emission flux. Gross photosynthesis follows a Michaelis--Menten response to atmospheric CO\textsubscript{2} with a temperature modifier,
\begin{equation}
P = k_p \, C_{v(s)} \, k_{MM}\, \frac{p\text{CO}_2 - k_c}{K_M + p\text{CO}_2 - k_c} \cdot \frac{(15 + \Delta T)^2 (25 - \Delta T)}{5625},
\qquad C_{v(s)} = C_{v,0} - k_D \int_{t_0}^{t} D\,dt',
\label{eq:photosynthesis}
\end{equation}
where $p\text{CO}_2$ is the atmospheric mixing ratio (ppmv), $k_c$ the compensation point, $K_M$ the half-saturation constant, $k_{MM}$ normalises pre-industrial photosynthesis, and $k_D$ is the fraction of land-use biomass loss that is permanent. Plant respiration follows an Arrhenius response, $R_P = k_r C_v k_A e^{-E_a/(R\,T_s)}$, and soil respiration the Lloyd--Taylor function, $R_S = k_{sr} C_s k_B\, e^{-308.56/(T_s - 227.13)}$, with $T_s = 288.15\,\text{K} + \Delta T$.

\subsubsection{Ocean}
Each ocean box carries three tracers: dissolved inorganic carbon (DIC), total alkalinity (TA) and nitrate (NO\textsubscript{3}). All tracers are advected by the water fluxes of \citet[Fig.~1, Table~2]{lentonLandOceanCarbon2000}: thermohaline overturning $k_T$ (intermediate $\to$ warm surface $\to$ cold surface $\to$ deep $\to$ intermediate), high-latitude overturning $k_U$ (intermediate $\to$ cold surface $\to$ deep $\to$ intermediate), and two-way exchange between warm surface and intermediate water ($k_{wi}$) and between intermediate and deep water ($k_{id}$). For a tracer with concentration $c_j$ in box $j$, the transport from box $j$ to box $k$ is $\rho\, Q_{jk}\, c_j$, where $Q_{jk}$ is the water flux and $\rho$ is seawater density.

Biological export follows \citet{knoxChangesAtmosphericCO21984}, whose ocean parameters Lenton adopts. All nitrate supplied to the warm surface box is consumed, and a fraction $\alpha_N = 0.01$ of the nitrate entering the cold surface box is consumed. Organic carbon is exported with C:N $= 7$, accompanied by CaCO\textsubscript{3} at a CaCO\textsubscript{3}:C\textsubscript{org} ratio of 0.25. Of the exported organic matter, 30\% is remineralised in intermediate water and 70\% in deep water; CaCO\textsubscript{3} dissolves entirely in deep water. Alkalinity changes by $-2$ per mole of CaCO\textsubscript{3} formed and $+1$ per mole of nitrate consumed, with the reverse on remineralisation. Ocean-mean concentrations are DIC $= 2250$, NO\textsubscript{3} $= 33$ and TA $= 2326\ \mu$mol\,kg$^{-1}$; the alkalinity inventory is set so that the model's pre-industrial steady state has an atmospheric mixing ratio of 280 ppmv, as in \citet{lentonLandOceanCarbon2000}.

Air--sea exchange in each surface box $j \in \{l, h\}$ is
\begin{equation}
F_j = k_s\, A_j\, \rho\, K_0(T_j)\left(p\text{CO}_{2,a} - p\text{CO}_{2,j}\right),
\label{eq:oceanflux}
\end{equation}
where $k_s$ is the piston velocity, $A_j$ the box surface area, $K_0$ the CO\textsubscript{2} solubility \citep{weissCarbonDioxideWater1974}, and $T_l = 295.15$\,K and $T_h = 275.15$\,K. Surface-water $p\text{CO}_{2,j}$ is computed from DIC and TA with the carbonate system, using the dissociation constants of \citet{goyetNewDeterminationCarbonic1989} and the borate constant of \citet{dicksonThermodynamicsDissociationBoric1990} at salinity 34.7, solving iteratively for [H\textsuperscript{+}]:
\begin{equation}
\text{TA} = \text{DIC}\,\frac{K_1[\text{H}^+] + 2K_1K_2}{[\text{H}^+]^2 + K_1[\text{H}^+] + K_1K_2} + \frac{B_T K_B}{K_B + [\text{H}^+]}, \qquad
p\text{CO}_{2,j} = \frac{1}{K_0}\,\frac{\text{DIC}\,[\text{H}^+]^2}{[\text{H}^+]^2 + K_1[\text{H}^+] + K_1K_2}.
\label{eq:carbonate}
\end{equation}

The atmosphere integrates the land and ocean fluxes and all emissions:
\begin{equation}
\frac{dC_a}{dt} = -P + R_P + R_S - F_l - F_h + E_{\text{ind}} + D(t) + J_{\text{pf}},
\label{eq:atm}
\end{equation}
where $E_{\text{ind}}$ is industrial CO\textsubscript{2}, $D(t)$ land-use CO\textsubscript{2} (both from DICE~2023, net of abatement), and $J_{\text{pf}}$ the permafrost flux (Eq.~\ref{eq:pf_resp}).

\subsubsection{Calibration and initial state}
We use Lenton's calibration of the CO\textsubscript{2} fertilisation parameter ($K_M = 120$ ppmv) and permanent land-use loss ($k_D = 0.27$). The model is started from its pre-industrial steady state in 1800 and forced with the historical emissions of \citet[Table~3]{lentonLandOceanCarbon2000} to 1990 and Global Carbon Budget fossil and land-use emissions from 1990 to 2019 \citep{friedlingsteinGlobalCarbonBudget2025}, with prescribed historical warming. The resulting 2020 state (atmospheric CO\textsubscript{2} 419 ppmv) initialises all carbon pools for the coupled runs, so that land and ocean reservoirs begin in a state consistent with their emissions history.

\subsubsection{Model evaluation}
With this configuration, the model reproduces the pre-industrial steady state of \citet[Fig.~1]{lentonLandOceanCarbon2000}: ocean carbon pools of 716, 139, 9{,}997 and 26{,}888 GtC (Lenton: 730, 140, 10{,}040 and 26{,}830 GtC), and gross air--sea fluxes of 56.9/58.4 GtC\,yr$^{-1}$ (warm, in/out; Lenton 57.0/57.9) and 19.0/17.5 GtC\,yr$^{-1}$ (cold; Lenton 19.1/18.2). Forced with Lenton's extended business-as-usual scenario (IS92a to 2100, declining linearly to zero by 2200) and coupled to his grey-atmosphere energy balance, it gives:
\begin{itemize}
\item 352.7 ppmv in 1990 (Lenton: 353), with land and ocean sinks of 1.98 and 1.99 GtC\,yr$^{-1}$ (1.98 and 1.98);
\item a terrestrial switch from carbon sink to source in 2113 (${\sim}$2120), with 68 GtC lost while the land is a source (69);
\item peak CO\textsubscript{2} of 984 ppmv in 2174 (985 ppmv, 2170) and peak warming of 5.53\,K in 2184 (5.5\,K, 2180) \citep[Sect.~5.1]{lentonLandOceanCarbon2000}.
\end{itemize}
Under an idealised 1\%\,yr$^{-1}$ CO\textsubscript{2} increase to 4$\times$CO\textsubscript{2}, the model's carbon--concentration feedbacks are $\beta_L = 0.69$ and $\beta_O = 0.86$ PgC\,ppm$^{-1}$, and its land carbon--climate feedback is $\gamma_L = -32$ PgC\,K$^{-1}$. The corresponding CMIP6 values are $1.22 \pm 0.40$, $0.79 \pm 0.07$ and $-58 \pm 38$ \citep{aroraCarbonConcentrationCarbon2020}; the ocean carbon--climate feedback is zero by construction in this model version.

\subsection{Permafrost carbon}
Permafrost follows the Hector v2.3pf implementation \citep{woodardPermafrostImplementationSimple2021, linkHectorPermafrostImplementation2021}. Frozen permafrost carbon ($C_{\text{perm}}$) thaws into a thawed pool ($C_{\text{thawed}}$) as the frozen fraction declines with high-latitude warming, $T_{HL} = w_f\,\Delta T$ with $w_f = 2$:
\begin{equation}
f_{\text{frozen}} = 1 - \Phi\!\left(\frac{\ln T_{HL} - \mu}{\sigma}\right), \qquad
\frac{dC_{\text{perm}}}{dt} = \frac{d f_{\text{frozen}}}{dt}\, C_{\text{perm}} - J_{\text{abrupt}},
\label{eq:frozenfrac}
\end{equation}
with $\mu = 1.67$ and $\sigma = 0.986$. Refreezing returns carbon from the thawed pool only, up to the amount available. Thawed carbon is respired as
\begin{equation}
J_{\text{pf}} = r\,(1 - f_{\text{static}})\, C_{\text{thawed}}\, Q_{10}^{\,T_{200}/10},
\label{eq:pf_resp}
\end{equation}
where:
\begin{itemize}
\item $r = 0.02$\,yr$^{-1}$ is Hector's base respiration rate;
\item $f_{\text{static}} = 0.74$ is the non-labile fraction;
\item $Q_{10} = 2$;
\item $T_{200}$ is the 200-year running mean of $T_{HL}$, which in Hector is not allowed to decline;
\item a fraction $f_{CH_4} = 0.023$ of $J_{\text{pf}}$ is released as methane, which we add to the atmospheric carbon pool.
\end{itemize}
The permafrost pools in 2020 are obtained by running the module from 1800 to 2020 with the historical warming reconstruction used for the carbon-cycle spin-up, starting from Hector's pre-industrial near-surface (0--3\,m) permafrost carbon of 865 PgC with no thawed carbon \citep[Table~1]{woodardPermafrostImplementationSimple2021}. This gives 692 PgC of frozen and 150 PgC of thawed permafrost carbon in 2020, consistent with Hector's calibration to a modern frozen stock of ${\sim}730$ PgC in 2010 \citep[Table~2]{woodardPermafrostImplementationSimple2021}. Permafrost emissions before 2020 are part of the observed historical carbon budget to which the carbon cycles are initialised and are not added separately. In the permafrost-strength experiment (below), the spin-up uses the same respiration rate as the projection, so the thawed pool in 2020 is consistent with each strength.

Abrupt thaw is represented following the Armstrong McKay et al.\ assessment of an additional ${\sim}50\%$ of gradual-thaw emissions, a ${\sim}1.5$\textdegree C threshold (range 1.0--2.3\textdegree C) and a ${\sim}200$-year timescale \citep{armstrongmckayExceeding15degCGlobal2022}. To avoid a global discontinuity at the threshold, onset is smooth:
\begin{equation}
J_{\text{abrupt}} = \frac{1}{\tau_a}\max\!\left(0,\; 0.5\, w(\Delta T)\, G - A\right), \qquad
w(\Delta T) = \tfrac{1}{2}\left[1 + \tanh\!\left(\frac{\Delta T - 1.5}{0.5}\right)\right],
\label{eq:abruptthaw}
\end{equation}
where $G$ is cumulative gradual thaw, $A$ cumulative abrupt thaw, and $\tau_a = 200$ years. Abrupt thaw is irreversible: cooling halts it but does not return carbon to the frozen pool.

The permafrost module can be coupled to either carbon cycle. With the Lenton model, $J_{\text{pf}}$ enters the atmosphere (Eq.~\ref{eq:atm}). With DFAIR, it is added to the emissions partitioned among the four DFAIR reservoirs (Eq.~\ref{eq:dfair_reservoir}).

\subsection{Model configurations and coupling}
Four physical configurations share the DICE~2023 economy, two-box temperature model and radiative forcing (Eq.~\ref{eq:F}):
\begin{itemize}
\item DICE~2023 (DFAIR);
\item DFAIR with permafrost;
\item DICE-C (Lenton carbon cycle);
\item DICE-CP (Lenton carbon cycle with permafrost).
\end{itemize}
In the Lenton configurations, the Lenton atmospheric carbon replaces DFAIR's in the forcing equation. All configurations are integrated in continuous time from 2020 to 2300 with a fourth-order Runge--Kutta scheme ($\Delta t = 0.05$ yr), solving the carbonate system at every function evaluation. We end the simulations in 2300 because the DICE~2023 control schedule changes after 2300 (the upper bound on the emissions control rate, which permits net negative emissions, is reduced to 1), so later behaviour reflects the policy schedule rather than the physical system; welfare, $U_{\text{acc}}$, is accordingly accumulated over 2020--2300.

\subsection{Experimental design}
The emissions control rate and savings rate are held at the DICE~2023 cost--benefit optimal trajectories from GAMS \citep{barragePoliciesProjectionsSocial2024}. Differences between configurations therefore reflect the physical system, not re-optimised policy.

\subsubsection{Damage and discounting}
Each configuration is evaluated under the standard DICE~2023 damage coefficient ($a_2 = 0.003467$) and under the coefficient derived from \citet{howardMethodologyMattersCareful2025} ($a_2 = 0.01$, Eq.~\ref{eq:hs25_a2}). Each is also evaluated under standard DICE~2023 discounting and under zero time preference ($\rho = 0$, $\beta_{clim} = 0$, precautionary term retained). This gives sixteen runs (Table~\ref{tab:sensitivity_params}). Welfare is reported as the percentage change in $U_{\text{acc}}$ (Eq.~\ref{eq:uacc}) relative to DICE~2023 under the same damage and discounting assumptions.

\subsubsection{Social cost of carbon}
The social cost of carbon in year $t$ is computed by the pulse method:
\begin{equation}
\text{SCC}(t) = -1000\,\frac{\partial U_{\text{acc}}/\partial E(t)}{\partial U_{\text{acc}}/\partial C(t)},
\label{eq:scc}
\end{equation}
where $\partial U_{\text{acc}}/\partial E(t)$ is the change in $U_{\text{acc}}$(2300) from an additional 1 GtCO\textsubscript{2} emitted uniformly during year $t$, and $\partial U_{\text{acc}}/\partial C(t)$ is the marginal welfare value of US\$1 trillion of consumption in year $t$, $\int_t^{t+1} 1000\,\text{STP}\,c^{-\eta}\,dt'$. The SCC is in 2019 US\$ per tCO\textsubscript{2}, holding the control and savings trajectories fixed.

\subsubsection{Permafrost strength}
The strength of the permafrost carbon feedback is uncertain by more than an order of magnitude. AR6 assesses 18 (3.1--41) PgC released per \textdegree C of warming by 2100 \citep{canadellGlobalCarbonOther2021}, and estimates including abrupt thaw and fire reach ${\sim}63$ PgC\,\textdegree C$^{-1}$ \citep{schadelPermafrostWildfireCarbon2026}. To span this range, we scale the respiration rate in Eq.~\ref{eq:pf_resp} by a factor $m$ ($r = 0.02\,m$, with $m$ from 0.1 to 50) in both permafrost configurations.

For each $m$ we report:
\begin{itemize}
\item the realised feedback strength, defined as cumulative permafrost release over 2020--2100 divided by 2100 warming;
\item the change in peak warming;
\item the change in $U_{\text{acc}}$;
\item the change in the 2020 SCC, each relative to the same carbon cycle without permafrost.
\end{itemize}
Each quantity is computed in all four damage and discounting combinations.

\subsubsection{Feedback attribution}
To attribute the difference between DICE-CP and DICE~2023 to individual processes, we switch six feedbacks in the Lenton carbon cycle on and off in all $2^6 = 64$ combinations (Table~\ref{tab:feedbacks}), under standard damage and discounting:
\begin{itemize}
\item \textbf{CO\textsubscript{2} fertilisation:} disengaged by holding the atmospheric CO\textsubscript{2} seen by photosynthesis at its 2020 value.
\item \textbf{Photosynthesis, plant respiration and soil respiration temperature responses:} each disengaged by holding the temperature anomaly in that term at its 2020 value.
\item \textbf{Ocean carbonate buffering:} disengaged by holding surface-ocean $p\text{CO}_2$ at its 2020 values, so that ocean uptake responds linearly to the atmospheric CO\textsubscript{2} excess.
\item \textbf{Permafrost:} disengaged by removing the permafrost module.
\end{itemize}
Freezing each dependence at its 2020 value, rather than at pre-industrial values, keeps every configuration mass-consistent and avoids an initial imbalance. For each feedback we report the mean change in peak warming (2020--2300) between runs with the feedback on and off, across all combinations of the other five.
